\documentclass[11pt]{article}
\usepackage[margin=25mm]{geometry}
\usepackage{amsmath,amssymb}
\usepackage{booktabs}
\begin{document}

\title{I01 prevalidation 15 result:\\Dimension dependence of paired support transfer}
\date{2026-10-04}
\maketitle

\section*{Implementation crosscheck}
The generalized $n=3$ circuit reproduces the stage-14 transfer entries with maximum
error $1.43\times10^{-17}$. Thus the changed clock width and circuit construction
have not introduced a new $n=3$ mechanism.

\begin{center}
\begin{tabular}{lrrr}
\toprule
support class & $n=2$ & $n=3$ & $n=4$\\
\midrule
same support & 0.001635829 & 0.001442169 & 0.000591890\\
conjugate partner & $2.6\times10^{-17}$ & $5.4\times10^{-17}$ & $3.7\times10^{-16}$\\
different pair, same sign & $1.2\times10^{-17}$ & 0.000387617 & 0.000169526\\
different pair, opposite sign & 0.001401342 & 0.000697812 & 0.000320087\\
\bottomrule
\end{tabular}
\end{center}

\section*{Conclusion}
Conjugate-partner transfer is a numerical zero at all tested dimensions. Same-support
and different-pair opposite-sign transfer remain nonzero at $n=2,3,4$. The
different-pair same-sign class is nonzero for $n=3,4$ but exactly suppressed for
$n=2$. Therefore the $n=3$ support pattern is neither a unique artefact nor a
dimension-independent law; the number of momentum pairs changes the interference
selection structure.

\section*{Scope}
This narrows the claim appropriately: the toy circuit exhibits a dimension-dependent
reference-access phenomenon. It is not evidence that physical space dynamically
chooses three dimensions.

\section*{Reproduction}
Code: \texttt{../src/relational\_atlas\_i01/holdout/qiskit\_rank\_holdout15\_dimension\_transfer.py}.\\
Output: \texttt{../results/qiskit\_rank\_holdout15\_dimension\_transfer.json}.
\end{document}
